% Options for packages loaded elsewhere
\PassOptionsToPackage{unicode}{hyperref}
\PassOptionsToPackage{hyphens}{url}
%
\documentclass[
]{article}
\usepackage{lmodern}
\usepackage{amssymb,amsmath}
\usepackage{ifxetex,ifluatex}
\ifnum 0\ifxetex 1\fi\ifluatex 1\fi=0 % if pdftex
  \usepackage[T1]{fontenc}
  \usepackage[utf8]{inputenc}
  \usepackage{textcomp} % provide euro and other symbols
\else % if luatex or xetex
  \usepackage{unicode-math}
  \defaultfontfeatures{Scale=MatchLowercase}
  \defaultfontfeatures[\rmfamily]{Ligatures=TeX,Scale=1}
\fi
% Use upquote if available, for straight quotes in verbatim environments
\IfFileExists{upquote.sty}{\usepackage{upquote}}{}
\IfFileExists{microtype.sty}{% use microtype if available
  \usepackage[]{microtype}
  \UseMicrotypeSet[protrusion]{basicmath} % disable protrusion for tt fonts
}{}
\makeatletter
\@ifundefined{KOMAClassName}{% if non-KOMA class
  \IfFileExists{parskip.sty}{%
    \usepackage{parskip}
  }{% else
    \setlength{\parindent}{0pt}
    \setlength{\parskip}{6pt plus 2pt minus 1pt}}
}{% if KOMA class
  \KOMAoptions{parskip=half}}
\makeatother
\usepackage{xcolor}
\IfFileExists{xurl.sty}{\usepackage{xurl}}{} % add URL line breaks if available
\IfFileExists{bookmark.sty}{\usepackage{bookmark}}{\usepackage{hyperref}}
\hypersetup{
  hidelinks,
  pdfcreator={LaTeX via pandoc}}
\urlstyle{same} % disable monospaced font for URLs
\usepackage{longtable,booktabs}
% Correct order of tables after \paragraph or \subparagraph
\usepackage{etoolbox}
\makeatletter
\patchcmd\longtable{\par}{\if@noskipsec\mbox{}\fi\par}{}{}
\makeatother
% Allow footnotes in longtable head/foot
\IfFileExists{footnotehyper.sty}{\usepackage{footnotehyper}}{\usepackage{footnote}}
\makesavenoteenv{longtable}
\setlength{\emergencystretch}{3em} % prevent overfull lines
\providecommand{\tightlist}{%
  \setlength{\itemsep}{0pt}\setlength{\parskip}{0pt}}
\setcounter{secnumdepth}{-\maxdimen} % remove section numbering

\author{}
\date{}

\begin{document}

\hypertarget{dot-product-logits-in-gpt-the-mathematically-precise-sentence}{%
\section{Dot-product logits in GPT: the mathematically precise
sentence}\label{dot-product-logits-in-gpt-the-mathematically-precise-sentence}}

\hypertarget{export-to-pdf-vs-code}{%
\subsection{Export to PDF (VS Code)}\label{export-to-pdf-vs-code}}

\begin{itemize}
\tightlist
\item
  If you have a Markdown-to-PDF extension installed (e.g., Markdown
  PDF), use its export command on this file.
\item
  Otherwise: open the preview, then use your system print dialog and
  select \textbf{Print to PDF}.
\end{itemize}

\hypertarget{clean-restatement}{%
\subsection{Clean restatement}\label{clean-restatement}}

\textbf{Yes. Exactly. That sentence is mathematically precise.}

I'll restate it cleanly, then justify it from first principles.

\pagebreak

\hypertarget{core-equation-one-line}{%
\subsection{Core equation (one line)}\label{core-equation-one-line}}

For token (i):

\[
\text{logit}_i = h^\top E_i
\]

Where:

\begin{itemize}
\tightlist
\item
  (h) = compressed context + analysis (final hidden vector)
\item
  (E\_i) = representation (embedding) of token (i)
\end{itemize}

This is a \textbf{dot product similarity}.

\begin{center}\rule{0.5\linewidth}{0.5pt}\end{center}

\hypertarget{why-this-works-science}{%
\subsection{Why this works (science)}\label{why-this-works-science}}

\hypertarget{what-h-really-is}{%
\subsubsection{1) What (h) really is}\label{what-h-really-is}}

(h) is \textbf{not} a summary sentence.

It is a \textbf{direction in embedding space} that encodes:

\begin{itemize}
\tightlist
\item
  semantics
\item
  syntax
\item
  positional constraints
\item
  long-range dependencies
\end{itemize}

It answers:

\begin{quote}
What kind of token should come next?
\end{quote}

\begin{center}\rule{0.5\linewidth}{0.5pt}\end{center}

\hypertarget{what-e_i-really-is}{%
\subsubsection{2) What (E\_i) really is}\label{what-e_i-really-is}}

Each token embedding (E\_i) is also a direction, learned so that:

\begin{itemize}
\tightlist
\item
  tokens used in similar contexts point similarly
\item
  grammar and meaning are embedded geometrically
\end{itemize}

\begin{center}\rule{0.5\linewidth}{0.5pt}\end{center}

\hypertarget{dot-product-meaning}{%
\subsubsection{3) Dot product meaning}\label{dot-product-meaning}}

Dot product measures:

\begin{itemize}
\tightlist
\item
  alignment
\item
  agreement
\item
  compatibility
\end{itemize}

\begin{longtable}[]{@{}ll@{}}
\toprule
Dot product value & Interpretation\tabularnewline
\midrule
\endhead
Large positive & Token fits context\tabularnewline
Near zero & Token is irrelevant\tabularnewline
Negative & Token contradicts context\tabularnewline
\bottomrule
\end{longtable}

Then \textbf{softmax} turns these compatibilities into probabilities.

\begin{center}\rule{0.5\linewidth}{0.5pt}\end{center}

\hypertarget{why-compressed-is-the-right-word}{%
\subsection{Why ``compressed'' is the right
word}\label{why-compressed-is-the-right-word}}

The entire history:

\[
(x_1, x_2, \dots, x_T)
\]

is mapped to:

\[
h \in \mathbb{R}^{500}
\]

This mapping is:

\begin{itemize}
\tightlist
\item
  lossy
\item
  task-oriented
\item
  predictive (not reconstructive)
\end{itemize}

Only information useful for predicting the next token survives.

\begin{center}\rule{0.5\linewidth}{0.5pt}\end{center}

\hypertarget{important-nuance}{%
\subsection{Important nuance}\label{important-nuance}}

This is \textbf{not} cosine similarity:

\begin{itemize}
\tightlist
\item
  norms matter
\item
  magnitude encodes confidence-like information
\end{itemize}

Temperature (and other sampling controls) later rescales this.

\begin{center}\rule{0.5\linewidth}{0.5pt}\end{center}

\hypertarget{one-line-summary}{%
\subsection{One-line summary}\label{one-line-summary}}

GPT predicts the next token by choosing the token whose embedding has
the highest dot-product alignment with the compressed, causally
accumulated context vector.

That's the model in one sentence.

\begin{center}\rule{0.5\linewidth}{0.5pt}\end{center}

\hypertarget{optional-next-topics}{%
\subsection{Optional next topics}\label{optional-next-topics}}

\begin{itemize}
\tightlist
\item
  why this equals nearest-neighbor under constraints
\item
  how attention learns to shape (h)
\item
  why magnitude is not literally confidence, but often correlates
\item
  why this beats symbolic rules
\end{itemize}

\end{document}
